\documentclass[11pt,openany,a4paper]{memoir}
\usepackage[T1]{fontenc}
\usepackage[utf8]{inputenc}
\usepackage{mathpazo}
\usepackage{microtype}
\newcommand{\olpath}{../upstream}
\input{\olpath/sty/open-logic.sty}
% Javanese text layer; Babel's English backend is used only for package compatibility.
% All displayed captions and token forms used in this tranche are overridden.
\setlocalecaption{english}{theorem}{Teorema}
\setlocalecaption{english}{example}{Tuladha}
\setlocalecaption{english}{definition}{Definisi}
\setlocalecaption{english}{lemma}{Lema}
\setlocalecaption{english}{proposition}{Proposisi}
\setlocalecaption{english}{corollary}{Korolari}
\setlocalecaption{english}{problem}{Soal}
\setlocalecaption{english}{problems}{Soal}
\setlocalecaption{english}{remark}{Cathetan}
\setlocalecaption{english}{note}{Cathetan}
\setlocalecaption{english}{axiom}{Aksioma}
\setlocalecaption{english}{case}{Kasus}
\setlocalecaption{english}{convention}{Konvensi}
\settexttoken{element}{anggota}{anggota}[Anggota][Anggota]
\definetoken{a}{element}{salah siji}
\definetoken{A}{element}{Salah siji}
\settexttoken{formula}{formula}{formula}[Formula][Formula]
\settexttoken{derivation}{derivasi}{derivasi}[Derivasi][Derivasi]
\definetoken{a}{formula}{salah siji}
\definetoken{A}{formula}{Salah siji}
\definetoken{a}{derivation}{salah siji}
\definetoken{A}{derivation}{Salah siji}
\settexttoken{surjective}{surjektif}{surjektif}[Surjektif][Surjektif]
\settexttoken{injective}{injektif}{injektif}[Injektif][Injektif]
\settexttoken{bijective}{bijektif}{bijektif}[Bijektif][Bijektif]
\definetoken{a}{surjective}{}
\definetoken{A}{surjective}{}
\definetoken{a}{injective}{}
\definetoken{A}{injective}{}
\definetoken{a}{bijective}{}
\definetoken{A}{bijective}{}
\settexttoken{surjection}{surjeksi}{surjeksi}[Surjeksi][Surjeksi]
\settexttoken{injection}{injeksi}{injeksi}[Injeksi][Injeksi]
\settexttoken{bijection}{bijeksi}{bijeksi}[Bijeksi][Bijeksi]
\definetoken{a}{surjection}{salah siji}
\definetoken{A}{surjection}{Salah siji}
\definetoken{a}{injection}{salah siji}
\definetoken{A}{injection}{Salah siji}
\definetoken{a}{bijection}{salah siji}
\definetoken{A}{bijection}{Salah siji}
\addto\captionsenglish{%
  \renewcommand{\contentsname}{Isine}
  \renewcommand{\chaptername}{Bab}
  \renewcommand{\partname}{Perangan}
  \renewcommand{\figurename}{Gambar}
  \renewcommand{\tablename}{Tabel}
  \renewcommand{\proofname}{Bukti}
  \renewcommand{\bibname}{Kapustakan}
}
\addto\extrasenglish{%
  \Crefname{thm}{Teorema}{Teorema}
  \Crefname{ex}{Tuladha}{Tuladha}
  \Crefname{defn}{Definisi}{Definisi}
  \Crefname{lem}{Lema}{Lema}
  \Crefname{prop}{Proposisi}{Proposisi}
  \Crefname{prob}{Soal}{Soal}
  \Crefname{rem}{Cathetan}{Cathetan}
  \Crefname{cor}{Korolari}{Korolari}
  \Crefname{axiom}{Aksioma}{Aksioma}
  \Crefname{note}{Cathetan}{Cathetan}
  \Crefname{case}{Kasus}{Kasus}
  \Crefname{conv}{Konvensi}{Konvensi}
  \Crefname{figure}{Gambar}{Gambar}
  \Crefname{table}{Tabel}{Tabel}
  \Crefname{section}{Perangan}{Perangan}
  \Crefname{chapter}{Bab}{Bab}
  \Crefname{page}{kaca}{kaca}
}
\newcommand{\JvReferenceName}[2]{\crefname{#1}{#2}{#2}\Crefname{#1}{#2}{#2}}
\AtBeginDocument{%
  \JvReferenceName{thm}{Teorema}
  \JvReferenceName{ex}{Tuladha}
  \JvReferenceName{defn}{Definisi}
  \JvReferenceName{lem}{Lema}
  \JvReferenceName{prop}{Proposisi}
  \JvReferenceName{prob}{Soal}
  \JvReferenceName{rem}{Cathetan}
  \JvReferenceName{cor}{Korolari}
  \JvReferenceName{axiom}{Aksioma}
  \JvReferenceName{note}{Cathetan}
  \JvReferenceName{case}{Kasus}
  \JvReferenceName{conv}{Konvensi}
  \JvReferenceName{figure}{Gambar}
  \JvReferenceName{table}{Tabel}
  \JvReferenceName{section}{Perangan}
  \JvReferenceName{chapter}{Bab}
  \JvReferenceName{page}{kaca}
  \renewcommand{\crefpairconjunction}{ lan }
  \renewcommand{\crefmiddleconjunction}{, }
  \renewcommand{\creflastconjunction}{, lan }
  \renewcommand{\crefrangeconjunction}{ nganti }
}

\setlrmarginsandblock{28mm}{28mm}{*}
\setulmarginsandblock{25mm}{28mm}{*}
\checkandfixthelayout
\linespread{1.05}
\hyphenpenalty=10000
\exhyphenpenalty=10000
\tolerance=2000
\emergencystretch=5em
% Keep each exercise at its original position using the default prob environment.
\pdfinfoomitdate=1
\pdftrailerid{}
\pdfsuppressptexinfo=15
\hypersetup{pdftitle={OpenLogic: Himpunan, Relasi, lan Fungsi - Basa Jawa},pdfauthor={Open Logic Project; Codex translation},pdflang={jv-Latn-ID}}
\begin{document}
\begin{titlingpage}
\raggedright
{\sffamily\Large OpenLogic \par}
\vspace{18mm}
{\sffamily\bfseries\Huge Himpunan, Relasi, lan Fungsi\par}
\vspace{6mm}
{\Large Basa Jawa, aksara Latin\par}
\vspace{18mm}
{\large Telung bab kapisan: dhasar himpunan, Paradoks Russell,
relasi, tatanan, graf, wit, lan fungsi.\par}
\vfill
Terjemahan mesin adhedhasar Open Logic Project.
Bab himpunan, relasi, lan fungsi wis diterjemahake kabeh; edhisi jangkep isih digarap.
Potret sumber rilis: 48 saka 722 berkas sumber wis diterjemahake.
Berkas sumber kang ndhukung pambangunan PDF iki: 26.
Unit sumber kang dimuat ing wacan PDF iki: 24 (pambuka lan telung bab).
\par\medskip
Sumber Inggris: OpenLogicProject/OpenLogic, revisi
\texttt{9620cc73f9c8e0ad003c514a5d3748f29611c4c0}.
\par\medskip
Karya asli lan terjemahan iki nganggo
\href{https://creativecommons.org/licenses/by/4.0/}{Creative Commons Attribution 4.0}.
Hak lan atribusi komponen asli tetep ditrapake.
Terjemahan iki ora nuduhake panyengkuyung saka Open Logic Project.
\end{titlingpage}
\pagestyle{plain}
\chapter*{Babagan Open Logic Project}
\addcontentsline{toc}{chapter}{Babagan Open Logic Project}


\textit{Open Logic Text} yaiku buku ajar sumber kabuka kang disusun
bebarengan, ngenani metalogika formal lan metodhe formal, wiwit tataran
madya (yaiku, sawise kuliah pambuka logika formal). Sanajan ditujokake
marang pamaca kang dhasar pasinaone dudu matematika (mligine mahasiswa
filsafat lan ilmu komputer), andharane tetep tliti.

Sawetara topik kang saiki wis dirembug bisa uga durung jangkep
andharane, lan akeh perangan isih mbutuhake revisi gedhe. Ing tembe,
kita ngrancang ngrembakakake teks iki supaya ngrembug luwih akeh topik.
Kita uga ngrancang nambahake glosarium, dhaptar wacan luwih lanjut,
cathetan sejarah, gambar, andharan kang luwih cetha, perangan kang
njlentrehake gayutane asil karo filsafat, ilmu komputer, lan matematika,
sarta luwih akeh soal lan tuladha ing teks iki.
Yen nemokake kaluputan utawa duwe usul,
\href{https://github.com/OpenLogicProject/OpenLogic/wiki/Contributing}{tulung ngandhani tim proyek}.

Proyek iki lumaku kanthi semangat sumber kabuka. Teks iki ora mung bisa
diakses kanthi bebas, nanging kita uga nyedhiyakake sumber LaTeX kanthi
Lisènsi Creative Commons Attribution. Lisènsi iki menehi hak marang
sapa wae kanggo ngundhuh, nggunakake, ngowahi, nata maneh, ngowahi
format, lan ngedum maneh karya kita, anggere menehi krédhit kang patut.
Katrangan liyane bisa diwaca ing situs Open Logic Project:
\href{http://openlogicproject.org/}{openlogicproject.org}.

\clearpage
\tableofcontents*
\olimport[../translation/content/sets-functions-relations/sets]{sets}
\olimport[../translation/content/sets-functions-relations/relations]{relations-complete}
\olimport[../translation/content/sets-functions-relations/functions]{functions}
\chapter*{Cathetan Sumber lan Notasi}
\addcontentsline{toc}{chapter}{Cathetan Sumber lan Notasi}

Ing tuladha relasi tatanan, sumber nganggo $I$ tanpa
ngenali jenenge luwih dhisik. Ing kono, $I$ kudu diwaca
minangka relasi identitas ing wilangan natural,
yaiku pasangan-pasangan ing diagonal kang wis diterangake.

Ing definisi pang wit, sumber nulis $z \in X \setminus B$,
nanging himpunan simpul wit wis dijenengi $A$ lan $X$
durung ditemtokake. Ing definisi iku, wacanen $X$ minangka $A$.
Maksimalitas banjur ateges saben simpul ing njaba $B$
ora bisa dibandhingake karo paling ora siji simpul ing $B$.

Notasi $R^+$ nduweni teges lokal kang beda:
ing perangan tatanan, tegese tutupan refleksif;
ing perangan operasi relasi, tegese tutupan transitif.
Gunakna definisi ing perangan kang lagi diwaca.

Andharan ekstensionalitas ing wiwitan kudu diwaca
bebarengan karo perangan Paradoks Russell:
ekstensionalitas njamin mung ana siji himpunan
kang cocog yen himpunan kasebut pancen ana.

Ing tuladha fungsi oyod kuadrat, sumber nyebut oyod kang dipilih
minangka positif, nanging domain wilangan natural kalebu nol.
Ing nol, wacanen iki minangka oyod kuadrat utama kang ora negatif.
Iki temuan audit sumber OLFUN-002 lan wis dibenerake ing teks terjemahan.
Ing tuladha himpunan pandhisik lan panerus, masukan dijenengi $n$,
nanging banjur ditulis $x$. Ing kono, $x$ nuduhake masukan $n$ kang padha.
Iki temuan OLFUN-003; teks terjemahan saiki nggunakake $x$ kanthi ajeg.

Proposisi yen saben injeksi duwe invers kiwa mbutuhake katrangan
tumrap domain kosong. Bukti ing sumber milih $a \in A$, mula
lumaku yen $A$ ora kosong. Yen $A$ lan $B$ loro-lorone kosong,
fungsi kosong dadi inverse dhewe. Nanging yen $A$ kosong lan $B$
ora kosong, injeksi kosong ora duwe invers kiwa saka $B$ menyang $A$.
Mula, wacanen proposisi iku kanthi sarat $A$ ora kosong utawa $B$ kosong.
Iki temuan OLFUN-001; awak proposisi terjemahan nganggo sarat prasaja
$A$ ora kosong lan cathetan jejer menehi sarat umum kang pas.
Panganggone Aksioma Pilihan ing bukti invers tengen wis diterangake
ing cathetan sikil sumber.

Ing andharan graf fungsi, sumber nyebut $R_f\subseteq A\times B$
minangka relasi "ing $A\times B$". Miturut definisi buku iki,
relasi ing himpunan kasebut malah dadi subhimpunan
$(A\times B)^2$. Teks terjemahan wis mbenerake iki dadi relasi
antarane $A$ lan $B$ kang kaemot ing $A\times B$ (OLFUN-004).

Watesan fungsi ing $C$ mung nyaring koordinat masukan, mula grafe
$R_f\cap(C\times B)$. Watesan relasi ing bagean sadurunge yaiku
$R\cap C^2$ lan nyaring loro koordinat. Rumus fungsi ing sumber bener;
pratelan yen loro gagasan persis padha wis diwatesi dadi analogi
kanthi bedane dijlentrehake (OLFUN-005).

Cathetan iki tambahan panyunting ing terjemahan mesin.
Rumus ing sumber selaras tetep dijaga supaya
bisa dibandhingake langsung karo sumber Inggris.

\bibliographystyle{natbib-oup}
\bibliography{open-logic}
\end{document}
